\subsection{Nutzerstudie - Aufbau/Tools}
Anhand der Nutzerstudie soll die erkennbarkeit von Halluzinierten Generierten antworten eines LLMs aufgezeigt werden. Zudem soll überprüft werden, welche bzw. ob ansätze zu einer steigerung der erkennbarkeit von Halluzinierten antworten führt und ob hierbei diese schneller vom nutzer erkannt werden könne. Zudem soll festgestellt werden, ob eine person sicherer beim erkennen von Halluzinationen wird, beim verwenden gewisser hilfsmittel? Und ob diese selbsteinschätzung mit der tatsächlichen einschätzung übereinstimmt.
\\
\subsubsection{Versuchsaufbau}
Die basis der Nutzerstudie bildet eine Shortstory und ein darauf basierender Fragenkatalog. Innerhalb dieses Fragenkatalogs werden nicht nur inhaltliche dinge abgefragt, sondern auch interpreationen aus dem bestehenden text. Bspw. "Emily fliegt nach London" => "Warum fliegt Emily nach Berlin?".
Diese Shortstory und der darauf basierende Fragenkatalog wird einem LLM zur verfügung gestellt, welches diese Fragen beantwortet. Hierbei geben wir dem LLM die anforderung, gegebene Antworten mit der entsprechenden textstelle zu referenzieren und zusätzlich den jeweiligen textabschnitt auszugeben und die wichtigen stellen die zur entscheidung geführt haben zu markieren.
Die Teilnehmer der Nutzerstudie werden nun in drei Gruppen aufgeteilt. Alle Gruppen erhalten hierbei die konkreten antworten welche das LLM gegeben hat und die jeweilige Shortstory. Zusätzlich gibt es die folgenden unterschiede innerhalb der drei Nutzergruppen:
\begin{enumerate}
    \item keine referenzierung
    \item Zeilen referenzierung
    \item Refernezierung mit markierter textstelle
\end{enumerate}
Die Aufgabe der Teilnehmer ist es jetzt, die Antworten des LLMs zu überprüfen und Halluzinierte Antworten zu identifizieren. Dabei Prüfen wir:
\begin{enumerate}
    \item Wie schnell die Teilnehmer die Antworten überprüfen können?
    \item Wie sicher sich die Teilnehmer fühlen beim identifizieren der Antworten?
    \item Wie viele Antworten die Teilnehmer korrekt einordnen?
\end{enumerate}







